% !TEX root = ../../../main.tex
% Block card: catalogue.event_windows -- "Event windows" (Part I, the blocks chapter).
\cardfield{What it does}
Cuts one window of fixed length around every event a detector found, lined up so the event
starts at the same point in every window.

\cardfield{When to use it}
After a detector (and ideally \blockref{interrogation.event_shape}), when you want the window
blocks (clustering, the classifier, window images, \blockref{catalogue.shape_families}) to work
on the events themselves instead of on a regular grid of windows.

\begin{plainterms}[Lining windows up on events]
Think of photographing runners at a finish line: each photo is taken a fixed moment before the
runner crosses, so in every photo the runner is in the same place. Here each event is the
runner. By default a fifth of the window comes before the event's onset and four fifths after,
so the onset sits at the same column in every window and the shapes can be laid over each
other.
\end{plainterms}

\cardfield{What goes in, what comes out}
In go the events (a span set). Out come windows, one per event that fits, each carrying the
event's measures (amplitude, width, recovery, if Event shape measured them) and a train,
validation or test role from a blocked split, as \blockref{preprocessing.sliding_windows}
makes it.

\cardfield{Settings that matter}
\emph{length\_s} is the window length in seconds. Left at 0 it is chosen for you: twice the
typical event duration, between one minute and two hours. That holds the onset, the trough and
about half the recovery, not the return to baseline; set it yourself when the whole recovery
matters. \emph{anchor} picks what is lined up: the onset, the extremum, or the span's start.
\emph{pre\_frac} is the share of the window before the anchor.

\cardfield{Watch out for}
Events too close to the start or end of the span to fit a whole window are left out, not
padded, and counted. Events closer together than one window give overlapping windows; when two
overlapping windows have different roles, that is a leak, counted on the block's page but not
refused.
